\section{The Optimal Quotes at NSE Parameters}
\label{sec:nse}

\subsection{Choosing the Risk Aversion}

The data pin down $\sigma$, $A$ and $k$; the risk aversion $\gamma$ is the dealer's choice. We set
$\gamma=\GammaRA$ per rupee by a simple rule: at the open, the two terms of the spread in
equation~\eqref{eq:spread} should be about equal for the liquid stock. For RELIANCE the
inventory-risk premium at the open is then \textrm{INR}~0.82 and the market-making edge is
\textrm{INR}~0.79, so neither consideration dominates the quotes. The same $\gamma$ is used for
KSCL, because risk aversion belongs to the dealer, not to the stock. Section~\ref{sec:mc} shows the
conclusions do not depend on this choice.

\subsection{What the Quotes Look Like}

Figure~\ref{fig:quotes} plots the bid and ask against inventory. Each share of inventory moves
both quotes by $\gamma\sigma^2(T-t)$, which at the open is \textrm{INR}~0.82 for RELIANCE and
\textrm{INR}~1.34 for KSCL. This is large relative to the half-spread, so the skew quickly runs into
the constraint $\delta\ge0$: a dealer long a few shares posts its ask at the mid, the most
aggressive price it is allowed, and moves its bid far away. The quotes are steeper at 09:15 than
at 12:15 because there is more time left in which the inventory could lose value.

\begin{figure}[!htbp]
\centering
\includegraphics[width=\textwidth]{model_quotes_vs_inventory}
\caption{Optimal bid and ask relative to the mid against inventory, at the open (solid, dark) and at
midday (light). The dashed line is the reservation price. When the dealer is long, the ask hits the
mid and stays there (the $\delta^a\ge 0$ constraint).}
\label{fig:quotes}
\end{figure}

Figure~\ref{fig:spreaddecomp} splits the total spread into its two parts across the session. The
edge term is constant; the inventory-risk premium falls linearly to zero at the close, so the dealer
quotes widest at the open and narrowest at the end. KSCL's spread is wider throughout: its book
thins out more slowly (lower $k$, larger edge) and its price is more volatile (larger premium). The
dashed line is the spread the market actually charges, measured by Roll's estimator. We return to the
gap between them in Section~\ref{sec:discussion}.

\begin{figure}[!htbp]
\centering
\includegraphics[width=\textwidth]{model_spread_decomposition}
\caption{Optimal total spread through the session, split into the market-making edge (light) and the
inventory-risk premium (dark), at the calibrated parameters and $\gamma=\GammaRA$. Dashed: the
market's average Roll spread.}
\label{fig:spreaddecomp}
\end{figure}

Figure~\ref{fig:heatmap} shows the full quoting rule as a map: the skew $r-s$ for every combination of
time and inventory. Red means the dealer prices the stock above the mid (it is short and wants to
buy); blue means below (it is long and wants to sell). The contours fan out to the right: the same
inventory produces less skew as the close approaches, because there is less time for the price to
move against it.

\begin{figure}[!htbp]
\centering
\includegraphics[width=\textwidth]{model_skew_heatmap}
\caption{Reservation price minus mid, $-q\gamma\sigma^2(T-t)$, over time of day and inventory.
Contours at $\pm0.5$, $\pm1$ and $\pm2$ rupees.}
\label{fig:heatmap}
\end{figure}
